\section{Development Environment}

\subsection{Preparation of the Working Environment}

Before starting the implementation, we prepared a local
development environment on Ubuntu. This allowed us to
install and test the necessary tools without consuming
our Azure credits during the initial development phase.

Python was selected as the main programming language,
with Flask as the web framework. Flask provides the
essential features required to build our application
while keeping its structure relatively simple.

We also configured the Azure CLI and Azure Functions
Core Tools to support the later integration and
deployment of cloud services.

\subsection{Development Tools}

Table~\ref{tab:environment} presents the main tools
available in our local development environment.

\begin{table}[htbp]
    \centering
    \caption{Development environment configuration}
    \label{tab:environment}

    \begin{tabular}{lll}
        \toprule
        \textbf{Technology} &
        \textbf{Version} &
        \textbf{Purpose} \\
        \midrule
        Ubuntu & 24.04.5 LTS & Operating system \\
        Python & 3.12.3 & Programming language \\
        Flask & 3.1.3 & Web framework \\
        Werkzeug & 3.1.9 & WSGI utilities \\
        Git & 2.43.0 & Version control \\
        VS Code & 1.146.0 & Development editor \\
        Azure CLI & 2.93.0 & Cloud management \\
        Functions Core Tools & 4.16.1 & Serverless development \\
        LaTeXmk & 4.83 & Report compilation \\
        \bottomrule
    \end{tabular}
\end{table}

\subsection{Python Virtual Environment}

A Python virtual environment was created for the project
to keep its dependencies separate from the packages
installed on the operating system.

After installing Flask, we recorded the installed
dependencies in a \texttt{requirements.txt} file.
This will allow both team members to reproduce the
same Python environment and will also be useful
when deploying the application to Azure.

\subsection{Version Control and Collaboration}

The project is managed using Git, with a remote repository
hosted on GitHub. This provides a common workspace for
the development team and allows us to track modifications
throughout the project.

A \texttt{.gitignore} file was configured to exclude
the virtual environment, sensitive configuration files
and locally generated files that do not need to be
included in the repository.

The application and the LaTeX report are maintained
within the same project directory so that documentation
can progress alongside implementation.